\documentclass[conference]{IEEEtran}
\IEEEoverridecommandlockouts

\usepackage{cite}
\usepackage{amsmath,amssymb,amsfonts}
\usepackage{graphicx}
\usepackage{textcomp}
\usepackage{xcolor}
\usepackage{booktabs}
\usepackage{multirow}
\usepackage{listings}
\usepackage[colorlinks=true,linkcolor=black,citecolor=black,urlcolor=black]{hyperref}

\lstdefinestyle{pycode}{
    language=Python,
    basicstyle=\ttfamily\scriptsize,
    breaklines=true,
    frame=single,
    rulecolor=\color{black!30},
    backgroundcolor=\color{black!3},
    xleftmargin=4pt, xrightmargin=4pt,
    aboveskip=4pt,   belowskip=4pt,
    keepspaces=true,
    showstringspaces=false,
    columns=fullflexible,
    keywordstyle=\color{blue!70},
    commentstyle=\color{green!50!black},
    stringstyle=\color{red!60!black}
}
\lstdefinestyle{reporttext}{
    basicstyle=\ttfamily\scriptsize,
    breaklines=true,
    frame=single,
    rulecolor=\color{black!30},
    backgroundcolor=\color{gray!6},
    xleftmargin=4pt, xrightmargin=4pt,
    aboveskip=4pt,   belowskip=4pt,
    keepspaces=true,
    showstringspaces=false,
    columns=fullflexible
}

\def\BibTeX{{\rm B\kern-.05em{\sc i\kern-.025em b}\kern-.08em
    T\kern-.1667em\lower.7ex\hbox{E}\kern-.125emX}}

\begin{document}

\title{Supplementary Material:\\
From Segmentation to Structured Reports:\\
Auditable Continual Learning for Automation in Medical Imaging Diagnostics}

\author{
\IEEEauthorblockN{Farhan Kashif, Ahmed Sultan, Muhammad Dawood Rizwan,
Muhammad Naseer Bajwa, Muhammad Wakil Shahzad, Muhammad Moazam Fraz}
}

\maketitle

\begin{abstract}
This document accompanies the main paper and provides: (A) the JSON Schema Draft-07 specification derived from the production schema file \texttt{schemas/tumor\_descriptor\_schema.json}; and (B) all five deterministic report templates reproduced directly from \texttt{src/report\_generation/templates.py}, together with the findings and impression generation logic from the same module. The complete source code is available in the project repository.
\end{abstract}

%=============================================================
\section*{Supplementary A: JSON Schema}
%=============================================================

\subsection*{A.1 Source File}

The schema is defined in \texttt{schemas/tumor\_descriptor\_schema.json} and is loaded at runtime by \texttt{src/json\_generation/schema\_validator.py}. It is validated against every descriptor produced by \texttt{TumorDescriptorGenerator} before downstream use.

\subsection*{A.2 Top-Level Structure}

Table~\ref{tab:schema_top} lists the five required top-level fields. Each corresponds to one builder method in \texttt{TumorDescriptorGenerator} (\texttt{src/json\_generation/descriptor\_generator.py}).

\begin{table}[htbp]
\caption{JSON Schema Top-Level Fields (all required)}
\begin{center}
\begin{tabular}{l l p{3.8cm}}
\toprule
\textbf{Field} & \textbf{Builder Method} & \textbf{Description} \\
\midrule
\texttt{patient\_info}       & \texttt{\_build\_patient\_info()}       & Case ID, age, sex, scan date \\
\texttt{imaging\_metadata}   & \texttt{\_build\_imaging\_metadata()}   & Modalities, scanner, resolution \\
\texttt{segmentation\_results} & \texttt{\_extract\_segmentation\_results()} & Volumes, confidence, centroids \\
\texttt{anatomical\_mapping} & \texttt{\_build\_anatomical\_mapping()} & Atlas regions, hemisphere, midline \\
\texttt{model\_metadata}     & \texttt{\_build\_model\_metadata()}     & Model name, version, timestamp \\
\bottomrule
\end{tabular}
\label{tab:schema_top}
\end{center}
\end{table}

\subsection*{A.3 BraTS Label Convention}

The segmentation label convention used in \texttt{\_extract\_segmentation\_results()} is:

\begin{itemize}
    \item Label \textbf{1}: Necrotic core / non-enhancing tumor core
    \item Label \textbf{2}: Peritumoral edema
    \item Label \textbf{4}: Enhancing tumor (ET)
    \item \textbf{Whole Tumor (WT)}: all labels $> 0$ = \{1, 2, 4\}
    \item \textbf{Tumor Core (TC)}: labels \{1, 4\}
\end{itemize}

\subsection*{A.4 Key Schema Definitions}

\subsubsection*{\texttt{tumor\_component} definition}
Each of the three tumor components (enhancing, necrotic, edema) follows this structure:

\begin{lstlisting}[style=pycode]
# From schemas/tumor_descriptor_schema.json -> definitions/tumor_component
# Populated by TumorDescriptorGenerator._analyze_component(seg_data, label, voxel_vol)
{
  "present":          bool,   # np.any(seg_data == label)
  "volume_mm3":       float,  # voxel_count * voxel_vol
  "voxel_count":      int,    # np.sum(mask)
  "confidence_score": float,  # mean sigmoid activation [0,1]
  "centroid_coords":  [x, y, z]  # np.argwhere(mask).mean(axis=0)
}
\end{lstlisting}

\subsubsection*{\texttt{affected\_region} definition}
Each atlas region overlap record (from \texttt{region\_analysis.py}, formatted by \texttt{\_format\_subregion\_mapping()}):

\begin{lstlisting}[style=pycode]
# From schemas/tumor_descriptor_schema.json -> definitions/affected_region
{
  "region_id":                  int,    # Harvard-Oxford numeric ID
  "region_name":                str,    # e.g., "Middle Frontal Gyrus"
  "percentage_involvement":     float,  # % of region covered by tumor mask
  "tumor_volume_in_region_mm3": float,  # absolute overlap volume
  "total_region_volume_mm3":    float,  # optional
  "hemisphere":                 str     # "left" | "right"
}
\end{lstlisting}

\subsubsection*{Clinical Features (optional, auto-derived)}
\begin{lstlisting}[style=pycode]
# Built by TumorDescriptorGenerator._extract_clinical_features()
{
  "mass_effect":                bool,  # total_tumor_volume_mm3 > 10000
  "ventricular_involvement":    bool,  # False (detection not yet implemented)
  "eloquent_area_involvement":  list,  # ["motor_cortex"] if "Precentral" in top-5 regions
  "estimated_grade":            str    # always "unknown" (requires clinical classifier)
}
\end{lstlisting}

\subsection*{A.5 Descriptor Generation Pipeline}

\begin{lstlisting}[style=pycode]
# src/json_generation/descriptor_generator.py -- generate_descriptor()
descriptor = TumorDescriptorGenerator(atlas_mapper=BrainAtlasMapper())
json_doc = descriptor.generate_descriptor(
    case_id      = 'BraTS-GLI-00006-100',
    seg_path     = 'segmentation.nii.gz',
    t1_path      = 't1.nii.gz',          # required for FLIRT registration
    modalities   = ['T1', 'T1ce', 'T2', 'FLAIR'],
    model_name   = 'MoME+',
    model_version= 'v1.0.0'
)
# Descriptor is automatically validated against
# schemas/tumor_descriptor_schema.json before return.
\end{lstlisting}

%=============================================================
\section*{Supplementary B: Report Templates}
%=============================================================

\subsection*{B.1 Architecture Overview}

Report generation is implemented in \texttt{src/report\_generation/} with three files:

\begin{itemize}
    \item \textbf{\texttt{templates.py}} — five report skeleton templates (\texttt{REPORT\_TEMPLATES[0--4]}) and two generation functions (\texttt{generate\_findings\_from\_json()}, \texttt{generate\_impression\_from\_json()})
    \item \textbf{\texttt{report\_generator.py}} — \texttt{TemplateBasedGenerator} (deterministic, zero-hallucination) and \texttt{ReportGenerator} (LLM-based with LoRA, used for NLG metric evaluation)
    \item \textbf{\texttt{dataset.py} / \texttt{metrics.py}} — BLEU/ROUGE/METEOR evaluation
\end{itemize}

The entry point is \texttt{TemplateBasedGenerator.generate(json\_data, template\_idx)} which calls \texttt{generate\_template\_report()} with \texttt{template\_idx \% 5} selecting the skeleton.

\subsection*{B.2 Generation Logic}

\subsubsection*{Findings Generation (\texttt{generate\_findings\_from\_json})}

The findings section is assembled as a list of deterministically chosen sentences. Each sentence uses \texttt{template\_variation \% 5} to index into a set of five phrasings per finding type (reproducing variation without free-text generation):

\begin{lstlisting}[style=pycode]
# src/report_generation/templates.py -- generate_findings_from_json()

tumor_descriptions = [          # indexed by var_idx (0-4)
    "A {type} tumor measuring {volume} cm3 is identified",
    "There is a {type} lesion measuring approximately {volume} cm3",
    "{type} mass lesion identified, measuring {volume} cm3",
    "An abnormal {type} mass measuring {volume} cm3 is present",
    "Focal {type} lesion measuring {volume} cm3"
]
location_descriptions = [
    "centered in the {location}",
    "primarily involving the {location}",
    "located in the {location}",
    "within the {location}",
    "affecting the {location}"
]
involvement_descriptions = [
    "with {overlap}% regional involvement",
    "demonstrating {overlap}% involvement of the region",
    "occupying approximately {overlap}% of the {location}",
    "with significant involvement ({overlap}%)",
    "involving {overlap}% of the {location}"
]
# JSON path used:
#   seg_results['enhancing_tumor']['volume_cm3']
#   anatomical['enhancing_tumor_regions'][0]['overlap_percent']
#   anatomical['enhancing_tumor_regions'][0]['name']
\end{lstlisting}

Subsequent sentences cover: tumor core volume, peritumoral edema ($WT\_vol - TC\_vol$), laterality, and mass effect flag.

\subsubsection*{Impression Generation (\texttt{generate\_impression\_from\_json})}

Five fixed impression strings (no free generation). The presence check uses:

\begin{lstlisting}[style=pycode]
has_tumor = any(
    seg_results.get(key, {}).get('volume_cm3', 0) > 0.1
    for key in ['enhancing_tumor', 'tumor_core', 'whole_tumor']
)
# If False -> "No evidence of intracranial mass lesion."
# If True  -> impressions[var_idx] (one of five fixed strings)
\end{lstlisting}

\subsection*{B.3 The Five Report Templates}

Templates are stored in \texttt{REPORT\_TEMPLATES} as a Python list. Each is a format string consuming three fields: \texttt{\{clinical\_history\}}, \texttt{\{modalities\}}, \texttt{\{findings\}}, \texttt{\{impression\}}. The \texttt{findings} and \texttt{impression} values are produced by the two generation functions described above; \texttt{clinical\_history} and \texttt{modalities} come directly from the JSON descriptor.

\subsubsection*{Template 0 (index 0): Standard Clinical Format}

\begin{lstlisting}[style=reporttext]
CLINICAL HISTORY: {clinical_history}

TECHNIQUE: Multi-sequence MRI brain examination including {modalities}.

FINDINGS:
{findings}

IMPRESSION:
{impression}
\end{lstlisting}

\subsubsection*{Template 1 (index 1): Detailed Structured Format}

\begin{lstlisting}[style=reporttext]
MRI BRAIN WITH AND WITHOUT CONTRAST

CLINICAL INDICATION: {clinical_history}

IMAGING TECHNIQUE: {modalities} sequences acquired.

FINDINGS:
{findings}

IMPRESSION:
{impression}
\end{lstlisting}

\subsubsection*{Template 2 (index 2): Concise Format}

\begin{lstlisting}[style=reporttext]
EXAMINATION: MRI Brain
INDICATION: {clinical_history}

FINDINGS:
{findings}

IMPRESSION:
{impression}
\end{lstlisting}

\subsubsection*{Template 3 (index 3): Comprehensive Format}

\begin{lstlisting}[style=reporttext]
BRAIN MRI REPORT

Clinical History: {clinical_history}

Technique: Multiparametric MRI including {modalities} sequences.

Findings:
{findings}

Impression:
{impression}
\end{lstlisting}

\subsubsection*{Template 4 (index 4): Academic / Conclusion Format}

\begin{lstlisting}[style=reporttext]
MRI BRAIN EXAMINATION

HISTORY: {clinical_history}

PROTOCOL: {modalities}

FINDINGS:
{findings}

CONCLUSION:
{impression}
\end{lstlisting}

\subsection*{B.4 Complete Generated Report Example (Template 0, Variation 0)}

The following is a representative output of \texttt{TemplateBasedGenerator.generate(json\_data, template\_idx=0)} for a right-sided high-grade glioma case:

\begin{lstlisting}[style=reporttext]
CLINICAL HISTORY: Brain tumor evaluation

TECHNIQUE: Multi-sequence MRI brain examination including T1, T1ce, T2, FLAIR.

FINDINGS:
A enhancing tumor measuring 8.5 cm3 is identified centered in the
Middle Frontal Gyrus with 42.3% regional involvement.
The tumor core measures 38.2 cm3.
Peritumoral edema extending into adjacent white matter measuring
approximately 34.2 cm3.
The lesion is right hemispheric in distribution.
No significant midline shift is observed.

IMPRESSION:
The imaging findings are consistent with a primary brain neoplasm,
likely high-grade glioma given the enhancement pattern and degree of
edema. Clinical correlation and tissue diagnosis recommended.
\end{lstlisting}

\noindent\textbf{Zero-hallucination guarantee:} Every numeric value in the FINDINGS section is retrieved directly from the validated JSON descriptor (volumes from \texttt{segmentation\_results}, region name and overlap from \texttt{anatomical\_mapping.subregion\_mapping.enhancing\_tumor[0]}). No language model inference is involved in the template pathway.

%=============================================================
\section*{Supplementary C: Companion File Index}
%=============================================================

\begin{table}[htbp]
\caption{Supplementary Files Provided}
\begin{center}
\begin{tabular}{l p{5cm}}
\toprule
\textbf{Filename} & \textbf{Contents} \\
\midrule
\texttt{supplementary\_A\_json\_schema.json} & Full JSON Schema Draft-07 (mirrors \texttt{schemas/tumor\_descriptor\_schema.json} with added annotations) \\
\texttt{supplementary\_B\_report\_templates.tex} & This document (LaTeX source) \\
\texttt{supplementary\_B\_report\_templates.pdf} & Compiled PDF \\
\midrule
\multicolumn{2}{l}{\textit{Source files (in repository):}} \\
\texttt{schemas/tumor\_descriptor\_schema.json} & Production schema \\
\texttt{src/json\_generation/descriptor\_generator.py} & JSON generation \\
\texttt{src/json\_generation/schema\_validator.py} & Runtime validation \\
\texttt{src/report\_generation/templates.py} & Template strings + generation logic \\
\texttt{src/report\_generation/report\_generator.py} & Generator classes \\
\texttt{src/models/gating\_network.py} & \texttt{HierarchicalGatingNetwork} + \texttt{ExpertFusion} \\
\texttt{src/atlas\_mapping/} & FLIRT registration + region analysis \\
\bottomrule
\end{tabular}
\end{center}
\end{table}

\end{document}
